\begin{frame}{How do particles move?}{A little bit of physics\dots}
    \begin{columns}[c]
        \begin{column}{0.55\textwidth}
            Consider a \textbf{particle} (coords.\ $\vec{x}$) moving in a \textbf{potential}
            \begin{equation*}
                U(\vec{x})
            \end{equation*}

            \includegraphics[height=12pt]{01_slides/img/brain.pdf}~Where will it land when it reaches a \textbf{stable equilibrium}?
        \end{column}
        \hfil
        \begin{column}{0.4\textwidth}
            \def\points{(0.1, 2.0) (0.5, 1.4) (0.9, 1.2) (1.3, 1.5) (1.7, 1.8) (2.1, 1.7) (2.5, 0.6) (2.9, 0.65) (3.3, 1.1) (3.7, 1.2) (4.1, 0.3) (4.5, 0.6)}
\begin{tikzpicture}[
        ax/.style={thick, CEAblack, ->}
    ]
    \draw[ax] (-0.2,0) -- (5.0,0) node[below] {$\vec{x}$};
    \draw[ax] (0,-0.2) -- (0,3.0) node[left] {$U$};
    \draw[thick, CEAblack, smooth] plot coordinates {\points};

    \draw<1>[thick, CEAred, fill=CEAred!15] (1.9, 1.95) circle (0.1);
    \draw<2->[thick, CEAred, fill=CEAred!15] (4.15, 0.42) circle (0.1);
\end{tikzpicture}

        \end{column}
    \end{columns}

    \vfill

    \begin{onlyenv}<2->
        The \textbf{equilibrium} is reached when
        \begin{equation*}
            \vec{x}^{*} = \arg\min_{\vec{x} \in \mathbb{R}^d} U(\vec{x})
            \quad \Rightarrow \quad
            \vec{F} = \vec{\nabla} U(\vec{x}) = \vec{0},
        \end{equation*}
        which gives us a good criterion to find the \textbf{optimum}.
    \end{onlyenv}
\end{frame}

\begin{frame}{How do particles move?}{A little bit of physics\dots}
    \begin{columns}[c]
        \begin{column}{0.55\textwidth}
            So now suppose you have a set of \enquote{parameters} $\Theta = \{ \theta_i \}_{i = 1}^d$, governed by a function
            \begin{equation*}
                L = L(\Theta)
            \end{equation*}
        \end{column}
        \hfil
        \begin{column}{0.4\textwidth}
            %\begin{tikzpicture}[
        ax/.style={thick, ->, CEAblack},
        ellipse/.style={thick, dashed},
        path/.style={very thick, ->, CEAred}
    ]
    % Axes
    \draw[ax] (-0.25, 0.0) -- (4.0, 0.0) node[below] {$\theta_1$};
    \draw[ax] (0.0, -0.25) -- (0.0, 4.0) node[left] {$\theta_2$};

    % Profile
    \node (opt) at (2.0, 2.0) {};
    \draw[ellipse, CEAblue, fill=CEAblue!15, rotate around={34:(opt)}] (opt) ellipse (1.0 and 2.0);
    \draw[ellipse, CEAgreen, fill=CEAgreen!15, rotate around={34:(opt)}] (opt) ellipse (0.7 and 1.4);
    \draw[ellipse, CEAred, fill=CEAred!15, rotate around={34:(opt)}] (opt) ellipse (0.4 and 0.8);

    % Optimum
    \node<2->[star, star points=5, star point ratio=2.5, draw, fill=yellow] at (opt) {};

    % Path
    \draw<2->[path] (4.1, 3.2) -- (3.0, 2.6) -- (2.9, 1.8) -- (2.5, 1.5) -- (2.2, 1.9);
\end{tikzpicture}

        \end{column}
    \end{columns}
\end{frame}
